% !PRINTED 520
\weberpage{520}
If the given functionals are actually fractions, that is, if they do not all reduce to integral functionals, then the common denominator must contain at least one prime factor that does not occur in all the numerators.

All this is in perfect agreement with the rules of elementary arithmetic, and the methods of proof that we have employed here are also essentially the same as those used there. The central point of the deduction is, on the one hand, the far-reaching generalisation of the concept of the unit, and on the other, the definition, founded upon it, of the greatest common divisor in \S.\,139.

\begin{center}
\S.\,142.\\
\gesperrt{Integral Functions in the Field $\overline{\Omega}$.}
\end{center}

The means we now have at our disposal suffice to apply to the functionals the mode of reasoning of which we made use in \S.\,2 of the first Band for the proof of \gesperrt{Gauss}'s theorem.

In fact, the theorem on which that proof essentially rests --- that a product of two integers is divisible by a prime only when this prime divides one of the factors --- has been proved valid also for the prime functionals now introduced.

We may accordingly, by following the proof just mentioned step by step, regard the theorem as proved:

1. \gesperrt{If two integral functions of a variable $t$, of degrees $h$ and $k$:}
\[
\varphi = \alpha_0 t^h + \alpha_1 t^{h-1} + \cdots + \alpha_h
\]
\[
\psi = \beta_0 t^k + \beta_1 t^{k-1} + \cdots + \beta_k,
\]
\gesperrt{whose coefficients are integral functionals that do not contain the variable $t$, have a product}
\[
\chi = \gamma_0 t^{h+k} + \gamma_1 t^{h+k-1} + \cdots + \gamma_{h+k}
\]
\gesperrt{in which the coefficients $\gamma_0, \gamma_1, \ldots, \gamma_{h+k}$ have a common prime divisor $\pi$, then $\pi$ must divide either all the coefficients of $\varphi$ or all the coefficients of $\psi$.}
